\subsection{Relation with affine Springer fibers}\label{ss:orb ASF}
From this subsection we restrict to the case $F$ is a local function field, i.e., $F=k\lr{t}$ for a finite field $k=\FF_{q}$. Let $\gamma\in\frg(F)$ be regular semisimple. The definitions of the affine Grassmannian and the affine Springer fiber $\sX_{\gamma}$ we gave in \S\ref{s:ASF} make sense when the base field $k$ is a finite field, so we have a sub-ind-scheme $\sX_{\gamma}$ of $\Gr_{G}$, both defined over $k$. 

The following lemma is clear from the definitions.
\begin{lemma}\label{l:X gamma k} The set of $k$-rational points $\sX_{\gamma}(k)$ is the same as the set $X_{\gamma}$ defined in \S\ref{sss:X gamma}, both as subsets of $\Gr_{G}(k)=G(F)/G(\calO_{F})$.
\end{lemma}

\subsubsection{} If we base change from $k$ to $\kbar$,  by Theorem \ref{th:ft quot} we know that $\L_{\gamma}\bs \sX_{\gamma,\kbar}$ is a proper algebraic space over $\kbar$. The proof there actually shows that this algebraic space is defined over $k$, which we denote by $\L_{\gamma}\bs \sX_{\gamma}$. We emphasize here that $\L_{\gamma}$ is viewed as an \'etale group scheme over $k$ whose $\kbar$-points is the plain group used to be denoted $\L_{\gamma}$.

In view of Lemma \ref{l:Ogamma} and Lemma \ref{l:X gamma k}, it is natural to expect that the orbital integral $O_{\gamma}(\one_{\frg(\calO_{F})})$ can be expressed as the number of $k$-points on the quotient $\L_{\gamma}\bs \sX_{\gamma}$. Such a relationship takes its cleanest form when $G=\GL_{n}$. 


\subsubsection{The case of $\GL_{n}$}\label{sss:GLn orb coho}
In the situation of \S\ref{sss:GLn orb formula}, upon choosing uniformizers $t_{i}\in F_{i}$, we defined the lattice $L_{0}=t^{\ZZ}_{1}\times \cdots\times t^{\ZZ}_{m}$. Base change from $k$ to $\kbar$, we may similarly define a lattice $\wt\L_{\gamma}$ as in \S\ref{sss:GLn centralizer}, using the same choice of uniformizers $t_{i}\in F_{i}$ (note that $F_{i}\wh\otimes_{k}\kbar$ may split into a product of fields, but $t_{i}\otimes 1$ will project to a uniformizer in each factor). The $\Gk$-action on $\wt\L_{\gamma}$ gives it the structure of an \'etale group scheme over $k$, just as $\L_{\gamma}$. We have $\wt\L_{\gamma}(k)=L_{0}$. There is an analog of Theorem \ref{th:ft quot} if we replace $\L_{\gamma}$ with $\wt\L_{\gamma}$. In particular, $\wt\L_{\gamma}\bs\sX_{\gamma}$ is a proper algebraic space over $k$ admitting a surjective map from a projective scheme.


\begin{theorem}\label{th:GLn orb} Let $G=\GL_{n}$. Let $\gamma$ be a regular semisimple element in $\frg(F)$. We fix the Haar measure on $G_{\gamma}(F)$ such that its maximal compact subgroup gets volume $1$. Then we have
\begin{equation}\label{GLn orb}
O_{\gamma}(\one_{\frg(\calO_{F})})=\#(\wt\L_{\gamma}\bs\sX_{\gamma})(k)=\sum_{i}(-1)^{i}\Tr\left(\Frob_{k}, \cohog{i}{\wt\L_{\gamma}\bs\sX_{\gamma,\kbar},\Ql}\right).
\end{equation}
\end{theorem}

The second equality in \eqref{GLn orb} follows from the Grothendieck-Lefschetz trace formula. Comparing \eqref{GLn orb} with \eqref{GLn L0}, we only need to argue that $L_{0}\bs X_{\gamma}=\wt\L_{\gamma}(k)\bs \sX_{\gamma}(k)$ is the same as $(\wt\L_{\gamma}\bs\sX_{\gamma})(k)$. Let us first make some general remarks about $k$-points on quotient stacks.

\subsubsection{$k$-points of a quotient}\label{sss:k pts quot} We consider  a quotient stack $\sY=[A\bs X]$ where $X$ is a scheme over $k$ and $A$ is an algebraic group over $k$ acting on $X$. Then, by definition, $\sY(k)$ is the groupoid of pairs $(S, f)$ where $S\to \Spec k$
is a left $A$-torsor, and $f:S\to X$ is an $A$-equivariant morphism. The isomorphism class of the $A$-torsor $S$ is classified by the Galois cohomology $\cohog{1}{k, A}:=\cohog{1}{\Gk,A(\kbar)}$. For each class $\xi\in \cohog{1}{k,A}$, let $S_{\xi}$ be an $A$-torsor over $k$ with class $\xi$. We may define a twisted form of $X$ over $k$ by $X_{\xi}:=A\bs (S_{\xi}\times X)$. We also have the inner form $A_{\xi}:=\Aut_{A}(S_{\xi})$ of $A$ acting on the $k$-scheme $X_{\xi}$. It is easy to see that $A$-equivariant morphisms $f:S_{\xi}\to X$ are in bijection with $X_{\xi}(k)$. Therefore we get a decomposition of groupoids
\begin{equation}\label{k pts quot}
\sY(k)=[A\bs X](k)\cong\bigsqcup_{\xi\in\cohog{1}{k,A}}A_{\xi}(k)\bs X_{\xi}(k).
\end{equation}


\subsubsection{Proof of Theorem \ref{th:GLn orb}} We have reduced to showing that $(\wt\L_{\gamma}\bs\sX_{\gamma})(k)= \wt\L_{\gamma}(k)\bs \sX_{\gamma}(k)$. By \eqref{k pts quot}, it suffices to show that $\cohog{1}{k,\wt\L_{\gamma}}=0$. We use the notation from \S\ref{sss:GLn orb coho}. From the definition of $\wt\L_{\gamma}$ we see that, as a $\Gk$-module, it is of the form
\begin{equation*}
\wt\L_{\gamma}\cong \bigoplus_{i=1}^{m}\Ind^{\Gk}_{\Gal(\kbar/k_{i})}\ZZ
\end{equation*}
where $k_{i}$ is the residue field of $F_{i}$. Therefore
\begin{equation*}
\cohog{1}{k,\wt\L_{\gamma}}\cong\bigoplus_{i=1}^{m}\cohog{1}{k_{i},\ZZ}=0.
\end{equation*}
\qed


\subsection{Stable orbital integrals} The setup is the same as \S\ref{ss:orb ASF}. For general $G$, the generalization of the formula \eqref{GLn orb} is not straightforward. Namely, the orbital integral $O_{\gamma}(\one_{\frg(\calO_{F})})$ by itself does not have a cohomological interpretation. The problem is that we may not be able to find an analog of $\wt\L_{\gamma}$ with vanishing first Galois cohomology so that $(\wt\L_{\gamma}\bs\sX_{\gamma})(k)$ is simply $\wt\L_{\gamma}(k)\bs \sX_{\gamma}(k)$. In view of formula \eqref{k pts quot},  a natural fix to this problem is to consider the twisted forms of $\sX_{\gamma}$ altogether. This suggests taking not just the orbital integral $O_{\gamma}(\one_{\frg(\calO_{F})})$ but a sum of several orbital integrals $O_{\gamma'}(\one_{\frg(\calO_{F})})$.


\subsubsection{Stable conjugacy} Fix a regular semisimple element $\gamma\in\frg(F)$. An element $\gamma'\in\frg(F)$ is called {\em stably conjugate} to $\gamma$ if it is in the same $G(\ov{F})$-orbit of $\gamma$. Equivalently, $\gamma'$ is stably conjugate to $\gamma$ if $\chi(\gamma')=\chi(\gamma)\in\frc(F)$. For $\gamma'$ stably conjugate to $\gamma$, one can attach a Galois cohomology class $\inv(\gamma,\gamma')\in\cohog{1}{F,G_{\gamma}}$ which becomes trivial in $\cohog{1}{F,G}$. The assignment $\gamma'\mapsto\inv(\gamma,\gamma')$ gives a bijection of pointed sets
\begin{equation}\label{stable class}
\{\gamma'\in\frg(F) \mbox{ stably conjugate to } \gamma\}/G(F) \cong \ker(\cohog{1}{F,G_{\gamma}}\to \cohog{1}{F,G}).
\end{equation}

\subsubsection{The case of $\GL_{n}$} When $G=\GL_{n}$, we use the notation from \S\ref{sss:GLn centralizer}. We have
\begin{equation*}
\cohog{1}{F,G_{\gamma}}\cong \prod_{i=1}^{m}\cohog{1}{F_{i},\Gm}=\jiao{1}.
\end{equation*}
Therefore, by \eqref{stable class}, all elements stably conjugate to $\gamma$ are in fact $G(F)$-conjugate to $\gamma$. Of course this statement can also be proved directly using companion matrices. 

\subsubsection{The case of $\SL_{2}$}\label{sss:st conj SL2} We consider the case $G=\SL_{2}$. Let  $a\in k^{\times}-(k^{\times})^{2}$. Let
\begin{equation*}
\gamma=\mat{0}{at}{t}{0}, \quad \gamma'=\mat{0}{at^{2}}{1}{0}
\end{equation*}
be two regular semisimple elements in $\sl_{2}(F)$. Since they have the same determinant $-at^{2}$, they are stably conjugate to each other. However, they are not conjugate to each other under $\SL_{2}(F)$. One can show that the stable conjugacy class of $\gamma$ consists of exactly two $\SL_{2}(F)$-orbits represented by $\gamma$ and $\gamma'$, see Exercise \ref{exIII:stable conj}.


\subsubsection{Definition of stable orbital integrals} Let $\varphi\in\sS(\frg(F))$. We define the {\em stable orbital integral} of $\varphi$ with respect to $\gamma$ to be
\begin{equation}\label{SO}
SO_{\gamma}(\varphi)=\sum_{\gamma'}O_{\gamma'}(\varphi)
\end{equation}
where $\gamma'$ runs over $G(F)$-orbits of elements $\gamma'\in \frg(F)$ that are stably conjugate to $\gamma$. For $\gamma'$ stably conjugate to $\gamma$, we have a {\em canonical isomorphism} $G_{\gamma}\cong G_{\gamma'}$ as $F$-groups. Therefore, once we fix a Haar measure on $G_{\gamma}(F)$, we get a canonical Haar measure on the other centralizers $G_{\gamma'}(F)$. It is with this choice that we define $O_{\gamma'}(\varphi)$ in  \eqref{SO}.

\subsubsection{Stable part of the cohomology}\label{sss:st H} The quotient group scheme $LG_{\gamma}/\L_{\gamma}$ acts on $\L_{\gamma}\bs\sX_{\gamma}$.  The component group $\pi_{0}(LG_{\gamma}/\L_{\gamma})$ is an \'etale group scheme over $k$ whose $\kbar$-points acts on $\cohog{*}{L_{\gamma}\bs\sX_{\gamma,\kbar},\Ql}$. We define the {\em stable part} of this cohomology group to be the invariants under this action
\begin{equation*}
\cohog{*}{\L_{\gamma}\bs\sX_{\gamma,\kbar}}_{\st}=\cohog{*}{\L_{\gamma}\bs\sX_{\gamma,\kbar}}^{\pi_{0}(LG_{\gamma}/\L_{\gamma})(\kbar)}.
\end{equation*}
It turns out that if we replace $\L_{\gamma}$ with any other $\Gk$-stable free abelian subgroup $\L\subset G_{\gamma}(F^{\ur})$ commensurable with $\L_{\gamma}$, the similarly defined stable part cohomology is canonically isomorphic to the above one.



\begin{theorem}[Special case of Goresky-Kottwitz-MacPherson {\cite[Th 15.8]{GKM}} and Ng\^{o} {\cite[Cor 8.2.10]{NgoFL}}]\label{th:orb coho} We have
\begin{equation}\label{d}
SO_{\gamma}(\one_{\frg(\calO_{F})})=\frac{1}{\vol(K_{\gamma},\mu_{G_{\gamma}})}\sum_{i}(-1)^{i}\Tr\left(\Frob_{k}, \cohog{i}{\L_{\gamma}\bs\sX_{\gamma,\kbar},\Ql}_{\st}\right).
\end{equation}
Here $K_{\gamma}\subset G_{\gamma}(F)$ is the parahoric subgroup of the torus $G_{\gamma}$. 
\end{theorem}

This cohomological interpretation of the stable orbital integral is the starting point of the proof of the Fundamental Lemma (see \cite{GKM} and \cite{NgoFL}). 


\subsubsection{}\label{sss:K gamma} Let us briefly comment on the definition of the parahoric subgroup $K_{\gamma}$. Let $\tilK_{\gamma}\subset G_{\gamma}(F)$ be the maximal compact subgroup. Then there is a canonical smooth group scheme $\calG_{\gamma}$ over $\calO_{F}$ whose $F$-fiber is $G_{\gamma}$ and whose $\calO_{F}$ points is $\tilK_{\gamma}$. This group scheme $\calG_{\gamma}$ is the finite type N\'eron model for the torus $G_{\gamma}$. Now let $\calG^{\circ}_{\gamma}\subset\calG_{\gamma}$ be the open subgroup scheme obtained by removing the non-neutral component of the special fiber of $\calG_{\gamma}$. Then, by definition, $K_{\gamma}=\calG^{\circ}_{\gamma}(\calO_{F})$. The positive loop group $L^{+}\calG^{\circ}_{\gamma}$ is, up to nilpotents, the neutral component of $LG_{\gamma}$.


\subsubsection{Sketch of proof of Theorem \ref{th:orb coho}}\label{sss:pf SO} We sketch an argument which is closer to that of \cite{GKM} than that of \cite{NgoFL}. One can show that there exists an \'etale $k$-subgroup $\wt\L\subset LG_{\gamma}$ commensurable with $\L_{\gamma}$, which maps {\em onto} the \'etale group scheme $\pi_{0}(LG_{\gamma})$ over $k$ (and the kernel is necessarily finite). The group $\wt\L$ is an analog of $\wt\L_{\gamma}$ defined in \S\ref{sss:GLn orb coho}. We form the quotient stack $[\wt\L\bs \sX_{\gamma}]$. Then we have an isomorphism
\begin{equation}\label{st coho quot}
\cohog{*}{\L_{\gamma}\bs\sX_{\gamma,\kbar},\Ql}_{\st}\cong \cohog{*}{[\wt\L\bs\sX_{\gamma,\kbar}],\Ql}.
\end{equation}
In fact, by the discussion in the end of \S\ref{sss:st H}, the left side above does not change if we replace $\L_{\gamma}$ by a commensurable lattice, so by shrinking $\L_{\gamma}$ we may assume $\L_{\gamma}\subset \wt\L$. On the other hand, the right side above can be computed by the Leray spectral sequence associated with the map $\L_{\gamma}\bs\sX_{\gamma,\kbar}\to[\wt\L\bs\sX_{\gamma,\kbar}]$ which is a torsor under the finite discrete group $(\wt\L/\L_{\gamma})(\kbar)$, therefore $\cohog{*}{[\wt\L\bs\sX_{\gamma,\kbar}],\Ql}=\cohog{*}{\L_{\gamma}\bs\sX_{\gamma,\kbar},\Ql}^{(\wt\L/\L_{\gamma})(\kbar)}$. Since $(\wt\L/\L_{\gamma})(\kbar)$ surjects onto $\pi_{0}(LG_{\gamma}/\L_{\gamma})(\kbar)$ by the choice of $\wt\L$, we see that
\begin{equation*}
\cohog{*}{\L_{\gamma}\bs\sX_{\gamma,\kbar},\Ql}^{\pi_{0}(LG_{\gamma}/\L_{\gamma})(\kbar)}=\cohog{*}{\L_{\gamma}\bs\sX_{\gamma,\kbar},\Ql}^{(\wt\L/\L_{\gamma})(\kbar)},
\end{equation*}
from which \eqref{st coho quot} follows.

By the discussion in \S\ref{sss:k pts quot} and \eqref{k pts quot}, we have
\begin{equation}\label{L bs X decomp}
[\wt\L\bs\sX_{\gamma,\kbar}](k)\cong\bigsqcup_{\xi\in\cohog{1}{k,\wt\L}} \wt\L(k)\bs\sX_{\gamma,\xi}(k).
\end{equation}
The inclusion $\wt\L\subset LG_{\gamma}$ gives
\begin{equation*}
\theta: \cohog{1}{k,\wt\L}\surj\cohog{1}{k,LG_{\gamma}}\cong\cohog{1}{F, G_{\gamma}}.
\end{equation*}
The first surjection follows from Lang's theorem because the quotient $LG_{\gamma}/\wt\L$ is connected; the second follows from another theorem of Lang which says that $\cohog{1}{F^{\ur}, G_{\gamma}}$ vanishes
\footnote{See \cite[Ch.X, \S7, p.170, Application and Example (b)]{Serre}. Let $K$ be a complete discrete valuation field with perfect residue field, and $K^{\ur}$ its maximal unramified extension. Then Lang's theorem asserts that $K^{\ur}$ is a $C_{1}$-field. Therefore $\cohog{1}{K^{\ur}, A}=0$ for any torus $A$ over $K$.}. For each $\xi\in\cohog{1}{k,\wt\L}$ such that $\sX_{\gamma,\xi}(k)\neq\varnothing$, one can show that the image of $\theta(\xi)$ in $\cohog{1}{F,G}$ is trivial. Therefore, by \eqref{stable class}, to each $\xi\in\cohog{1}{k,\wt\L}$ we can attach an element $\gamma'$ stably conjugate to $\gamma$, unique up to $G(F)$-conjugacy, such that $\inv(\gamma,\gamma')=\theta(\xi)$. One can show that $\sX_{\gamma,\xi}(k)$ is in bijection with the set $X_{\gamma'}$. Therefore,  \eqref{L bs X decomp} implies
\begin{equation}\label{pre SO}
\#[\wt\L\bs\sX_{\gamma}](k)=\#\ker(\theta)\sum_{\gamma'}\#(\wt\L(k)\bs X_{\gamma'})
\end{equation}
where the sum is over the $G(F)$-orbits of those $\gamma'$ stably conjugate to $\gamma$. Applying Lemma \ref{l:Ogamma} to the discrete cocompact subgroup $\wt\L(k)$, we have
\begin{equation*}
O_{\gamma'}(\one_{\frg(\calO_{F})})=\frac{1}{\vol(G_{\gamma}(F)/\wt\L(k),\mu_{G_{\gamma}})}\#(\wt\L(k)\bs X_{\gamma'}).
\end{equation*}
Plugging this into the right side of \eqref{pre SO}, we get
\begin{equation}\label{next SO}
\SO_{\gamma}(\one_{\frg(\calO_{F})})=\frac{1}{\#\ker(\theta)\vol(G_{\gamma}(F)/\wt\L(k),\mu_{G_{\gamma}})}\#[\wt\L\bs\sX_{\gamma}](k).
\end{equation} 
By the Grothendieck-Lefschetz trace formula,   $\#[\wt\L\bs\sX_{\gamma,\kbar}](k)$ is equal to the alternating Frobenius trace on $\cohog{*}{[\wt\L\bs\sX_{\gamma,\kbar}],\Ql}$, which, by \eqref{st coho quot}, can be identified with $\cohog{*}{\L_{\gamma}\bs\sX_{\gamma,\kbar},\Ql}_{\st}$. Therefore the theorem follows from the identity \eqref{next SO} together with the volume identity
\begin{equation}\label{vol id}
\#\ker(\theta)\vol(G_{\gamma}(F)/\wt\L(k),\mu_{G_{\gamma}})=\vol(K_{\gamma},\mu_{G_{\gamma}}).
\end{equation}
To show this, let $C=L^{+}\calG^{\circ}_{\gamma}\cap \wt\L$ (where $\calG^{\circ}_{\gamma}$ is the connected N\'eron model of $G_{\gamma}$ whose $\calO_{F}$ points is $K_{\gamma}$, see \S\ref{sss:K gamma}). This is a finite \'etale group over $k$. We have a short exact  sequence of group ind-schemes over $k$
\begin{equation*}
1\to C\to L^{+}\calG^{\circ}_{\gamma}\times\wt\L\to (LG_{\gamma})^{\red}\to 1
\end{equation*}
The associated six term exact sequence for $\Gk$-cohomology gives
\begin{equation*}
1\to C(k)\to K_{\gamma}\times \wt\L(k)\to G_{\gamma}(F)\to \cohog{1}{k,C}\to \ker(\theta)\to1
\end{equation*}
from which we get \eqref{vol id}, using that $\#C(k)=\#\cohog{1}{k,C}$. \qed



